\section{Discussion}
\label{sec:discussion}

\paragraph{Why organization would be phase-specific.} The two phases
demand different things. Producing a first answer is a \emph{generation}
problem: there is no incumbent candidate, and three agents arriving at the
same place by three routes costs nothing but tokens, since the pipeline
emits one answer and it is as good as the best route that reached it. A
division of labor has nothing to do here that redundancy does not already
do. Repairing a wrong answer is a \emph{differentiation} problem: an
incumbent candidate exists, it is in every agent's context, and the useful
move is to produce something unlike it. gemma-4's role-free team fails
there not for lack of effort --- it spends its whole budget --- but for
lack of difference: three agents that are now the same solver, all
conditioned on the same incumbent answer, reproduce it. A role contract
supplies a reason for one agent's next output not to be a function of the
previous agent's, worth nothing when there is nothing to differ from and a
great deal when there is; enforcement, the same prohibition moved to the
interface, moves the same quantity. The contract is necessary rather than
sufficient. With roles, Llama-3.3's Reviewer declines in $4.88$ of $5.93$
cycles, yet $0.750$ of its first-pass-wrong trajectories never write down a
second distinct answer and only $10$ of $276$ submit again: the role opens
the revision channel, and this solver has little to send through it.

\paragraph{What a role diagram is evidence of.} A diagram is a
hypothesis, and task-outcome evaluation --- the only evaluation most
multi-agent systems receive --- cannot test it: an aggregate score averages
a stage where the organization is idle with one where it does all its work,
dominated by whichever has more mass. Our system lost its entire
organization and moved the headline metric by $0.0338$. Two cheap additions
fix most of it: report performance \emph{conditioned on the first attempt
having failed}, computable from logs any retry-capable system already
writes, and \emph{measure} the realized organization rather than assume the
specified one --- indistinguishable by accuracy, trivially distinguishable
by whether all three agents are writing down answers.

\paragraph{What our own measurement errors imply.} Three instruments here
reported confidently on what they could not see: an exact-duplication
metric returned $0.000$ for actors a model judge found collapsed; a
logging defect reduced one actor's planner to $37$ visible characters,
which we nearly reported as a behavior; and the adjudication built to
check the first leaked the condition it was meant to be blind to. Each
produced a clean, plausible reading while not measuring the thing. A claim
that a system's structure is intact should therefore state what the
instrument would have to see to report otherwise --- and a blinding claim
should be tested on the bytes the judge saw, not the substitution applied
to them.
